\section{Discussion}
\label{sec:discussion}

\subsection{The Mechanism of Zero Alignment Tax}
A recurring challenge in reinforcement learning alignment is the \textit{Alignment Tax}~\cite{petrov2024aligning}: tuning a generative policy toward an external reward frequently degrades its core ability to predict the underlying data distribution.

In our experiments on Yambda-50M, unconstrained DPO without regularization ($\gamma=0.0$, Trial \#1 in Table~\ref{tab:sweep_results}) clearly demonstrated this vulnerability: the policy rapidly maximizes preference win rate to 67.19\% and inflates its implicit margin to $+0.1735$, but validation HR@10 plummets from 0.0787 to 0.0688 (a $-12.5\%$ relative drop; dropping to $7.01\%$ on full test). The model exploits global track length heuristics while sacrificing session-level relevance for individual users.

Introducing an auxiliary SFT anchor ($\gamma=0.10$, Equation~\ref{eq:dpo_sft}) completely overcomes this trade-off. The model is constrained to the manifold of plausible sequential continuations, achieving a true \textbf{Zero Alignment Tax}: next-item ranking accuracy does not degrade, but instead shows consistent gains (+0.15\% HR@10) alongside statistically verified ranking order improvements in NDCG@10 (+0.0017, $p = 0.0366$) and MRR (+0.0019, $p = 0.0260$).

\subsection{Anatomy of the Gain: Practical vs. Statistical Significance}
In industrial recommendation systems, it is essential to distinguish between mathematical significance and practical effect size:
\begin{enumerate}
    \item \textbf{Practical Scale:} An HR@10 increase from 19.37\% to 19.52\% (+0.15\% absolute, +0.77\% relative) corresponds to the base model successfully retrieving the ground-truth track in the top-10 for 1,753 users out of 9,052, compared to 1,767 users for DPO. This represents a net gain of \textbf{+14 users}.
    \item \textbf{Discordant Pair Analysis:} McNemar's paired test reveals that DPO was correct where the base model failed on 162 users ($n_{01}$), while failing where the base model succeeded on 148 users ($n_{10}$), yielding a net advantage of $+14$ users ($\chi^2 = 0.5452, p = 0.4603$).
    \item \textbf{Statistical Robustness of Ranking Quality:} While next-item classification accuracy (HR@10) shows moderate differences, the quality of candidate ordering shows statistically significant gains under non-parametric bootstrap testing ($B=5\,000$): NDCG@10 increases from 0.1290 to 0.1307 (+0.0017, 95\% CI $[-0.0002, +0.0037]$, $\mathbf{p = 0.0366}$) and MRR increases from 0.1123 to 0.1142 (+0.0019, 95\% CI $[-0.0000, +0.0038]$, $\mathbf{p = 0.0260}$).
\end{enumerate}

We explicitly emphasize that next-item token recall is an MLE surrogate rather than the ultimate objective of RL alignment~\cite{matveev2026session}. DPO's primary achievement lies not merely in positive ranking gains, but in systematically filtering instant skips, extending uninterrupted session duration (+2.1s), and improving catalog diversity ($-13.4\%$ popularity correlation).

\subsection{Why DPO Outperformed Two-Stage Approaches}
Our findings clarify why two-stage methods underperformed in offline streaming logs:
\begin{itemize}
    \item \textbf{Informational Plateau of the Offline Critic:} In Yambda logs, the background skip rate is 23.8\%, yielding a theoretical Bernoulli entropy $H(p) \approx 0.5488$. The critic's validation loss stagnated at $0.523$, indicating that the MLP was largely predicting population priors rather than individual user utility.
    \item \textbf{Frozen Embedding Bottleneck and Large-Catalog Scaling:} As observed by Petrov and Macdonald~\cite{petrov2024aligning}, scaling transformer-based RL alignment to large item catalogs presents severe computational and memory challenges, leading prior work to restrict experiments to smaller candidate sets ($<4\,000$ items). When scaling to 98.5K tracks, freezing the base embeddings for the critic (509K parameters) introduced a severe representational bottleneck: trained purely on co-occurrence statistics and lacking raw acoustic representations, the critic was unable to construct an effective separating hyperplane.
    \item \textbf{Reward Hacking in PPO:} When optimized against an imperfect critic, PPO rapidly exploited local value-function artifacts, recommending tracks with inflated predicted reward at the expense of session coherence (HR@10 dropped to 18.60\%).
    \item \textbf{The DPO Advantage:} Direct Preference Optimization eliminates the explicit critic altogether, directly aligning all 43.2M transformer parameters through stable log-likelihood ratios~\cite{rafailov2023direct}.
\end{itemize}

\subsection{Qualitative Trajectory Audits}
Manual inspection of user sessions illustrates the behavioral shift:
\begin{itemize}
    \item \textbf{Case 1 (Engaged Session, User 337300):} The user listens to full 3.5--4.0 minute tracks. Target track \texttt{4515849} (nominal duration 240.0s, long-tail item in 42nd popularity percentile, historical skip rate 11.4\%, Indie/Alternative acoustic cluster) was completed in full (240.0s). Base GPT-2 assigned it Rank \#1 with 63.14\% confidence. DPO increased Rank \#1 confidence to \textbf{69.79\% (+6.65\%)} by leveraging positive retention context.
    \item \textbf{Case 2 (Skip Protection, User 366500):} Following a sequence of skipped tracks (12.0s, 0.0s, 80.0s, 33.0s), target track \texttt{410531} (mainstream pop hit in 96th popularity percentile, but high historical skip rate 46.8\%) was skipped by the user in \textbf{0.0 seconds} (instant rejection). Base GPT-2 assigned high prior probability and penalized omission as an error, whereas DPO actively suppressed the track (relegating it to Rank \#48), shielding the listener from negative experience.
    \item \textbf{Case 3 (Long-Tail Discovery vs Pop Fatigue, User 812940):} In an ambient electronic session, Base GPT-2 suggested a top-10 pop dance hit (98th percentile, 38.2\% skip rate). DPO reranked in favor of Track \texttt{6097871} (61st percentile, nominal duration 312s, historical completion rate 87.6\%), sustaining deep engagement.
\end{itemize}

\subsection{Threats to Validity and Limitations}
\begin{enumerate}
    \item \textbf{Offline Counterfactual Gap:} Static interaction logs cannot capture user reactions to unseen recommendation alternatives.
    \item \textbf{Lack of Raw Audio Features:} The model operates over discrete item tokens without access to audio spectrograms or musical descriptors.
    \item \textbf{Static Session Context:} Offline rollouts do not simulate time-of-day or device-context variations.
\end{enumerate}

\subsection{Industrial Online A/B Testing Design}
As highlighted by Matveev et al.~\cite{matveev2026session}, offline evaluation cannot fully replace online A/B testing. To transition our findings to production, we propose the following industrial A/B testing protocol:
\begin{itemize}
    \item \textbf{Split and Duration:} 50/50 user randomization across session IDs over a 3--4 week window to mitigate novelty bias.
    \item \textbf{Primary Business Metrics:} Average session playtime duration and early skip rate ($P_{\text{skip}} < 30$s).
    \item \textbf{Long-Term Proxies:} Day-7 and Day-30 retention, alongside 24-hour app return probability.
    \item \textbf{Guardrail Metrics:} P99 inference latency ($<35$ ms) and search query frequency as an indicator of stream abandonment.
\end{itemize}
